% ============================================================
% 分片：s04-geometry  对应原文 PDF 第 44--48 页（印刷页码 44--48）
% 负责章节：4. SOME GEOMETRICAL BACKGROUND（几何背景）
% 译者：t08-s4   审校状态：初稿
% 审校：t06-s33 已完成逐式核对（修正 1 处）
% ============================================================
% 记法说明（不改符号，仅统一字形）：原文以粗体 $\mathbb{C}^n$、$\mathbb{C}^*$、
% $\mathbf{CP}^{n-1}$、$\mathbb{R}$ 排版，本译本按 GLOSSARY 与 SPEC 4.5 统一写作
% $\mathbb{C}^n$、$\mathbb{C}^*$、$\mathbf{CP}^{n-1}$、$\mathbb{R}$；指数、上下标、
% 权重、$D$ 与 $\tilde{D}$ 的符号一律照原文。原文明写直立 d 的微分形式，译文中
% 一律用 $\mathrm{d}$；原文中的 $i$ 为斜体，译文保持斜体 $i$。
% 第 4 节共 3 个带号公式：(4.1)--(4.3)，全部锁号。

% ==== 原文 PDF 第 44 页（印刷页 44）====
\section{几何背景}

\begin{keytakeaway}
\textbf{核心物理图像：}
第 3 节的构造可以理解为对 $\mathbb{C}^{n+1}$ 上紧群 $U(1)$ 的辛商（其对偶的全纯描述即对 $\mathbb{C}^*$ 取商）：先用 $D$ 项方程（矩映射取固定值 $r$）挑出一个截面，再对紧规范群 $U(1)$ 取商。不同的 $r$ 给出拓扑不同、却双有理等价的靶空间，因此线性 $\sigma$ 模型能在这些靶空间之间插值。
\par\textbf{数学工具：}
辛商（Marsden--Weinstein 约化）、几何不变量理论（GIT）、Kempf--Ness 定理、Duistermaat--Heckman 定理、双有理等价与胀开（blow-up）。
\par\textbf{主要结论：}
$r>0$ 时商空间是线丛 $\mathcal{O}(-n)$ 的全空间，$r<0$ 时是 $\mathbb{C}^n/\mathbb{Z}_n$，二者均为非紧 Calabi--Yau 流形；$r=0$ 时商奇异。三者在一个稠密开集上重合，且 $\mathbb{C}^*$ 不变的全纯函数（超势）可下降到每一个商上。这正是第 3 节插值手法、以及第 5.5 节时空拓扑改变的几何基础。
\end{keytakeaway}

第 3 节的构造有许多推广，其中一些将在第 5 节中探讨。不过，作者首先想解释一些几何背景，它们能为这一主题提供某些启示。这些事实对于阅读第 5 节并非严格必需，但澄清了其中若干问题，尤其是当我们着手考虑第 5.5 节中的拓扑改变时。实际上，我们将以一种初等的、\emph{ad hoc} 的方式，说明（见 [33] 或 [34, p.~158]）辛商（symplectic quotient）与全纯商（holomorphic quotient）之间的关系。

考虑流形 $Y=\mathbb{C}^{n+1}$，其坐标为 $s_1,\ldots,s_n$ 以及 $p$，并带有 $\mathbb{C}^*$ 作用
\begin{equation}
\begin{aligned} s_i &\to \lambda s_i\\ p &\to \lambda^{-n}p \end{aligned}
\tag{4.1}
\end{equation}
其中 $\lambda\in\mathbb{C}^*$。我们希望构造 $Y$ 对 $\mathbb{C}^*$ 的商。这样做的相关性将逐渐显现。既然 $\mathbb{C}^*$ 在 $Y$ 去掉原点 $O$ 上自由作用，人们或许会以为可以直接构造一个合理的商 $(Y-O)/\mathbb{C}^*$。对于紧群的自由作用，合理的商总是存在的；但对非紧群，情形则大不相同。

% [Tutor Note: (4.1) 的权重为 $Q_s=1$（$n$ 个坐标）与 $Q_p=-n$，故权重和 $\sum_i Q_i=n\cdot1+(-n)=0$。这一“权重和为零”既是第 47 页 $(n+1)$ 形式 $\Theta$ 的 $\mathbb{C}^*$ 不变性的来源，也是 $r<0$ 分支残余 $\mathbb{Z}_n\subset SU(n)$ 的来源。另注意 $\mathbb{C}^*$ 非紧：$|\lambda|$ 可取任意正实数，轨道在 $Y$ 中不闭，这正是下面集合论商非豪斯多夫的根源。]

% [原文疑误: 原文把 symplectic 两处误拼为 "sympectic"——第 46 页 "the sympectic structure of $Y//U(1)$" 与第 47 页 (ii) 中 "Its sympectic quotient"；而第 46 页首段 "the symplectic quotient" 等处拼写正确，故判定为排版笔误，译文统一按“辛”处理。]

设 $P$ 为 $Y$ 中满足 $p=0$ 的点，$P'$ 为 $Y$ 中满足 $s_1=\ldots=s_n=0$ 的点。在 $\mathbb{C}^*$ 作用下，$P$ 可被任意接近地移到原点附近（取 $\lambda\to0$），
% ==== 原文 PDF 第 45 页（印刷页 45）====
\label{op:45}
% [跨页接续: 原文第 44 页/第 45 页；本句跨页，此处按原文不切分句子]
对 $P'$ 也可同样为之（取 $\lambda\to\infty$）。即便在构造商时挖去原点，$P$ 的任一邻域中仍有一条 $\mathbb{C}^*$ 轨道与 $P'$ 的任一邻域相交。这意味着集合论商 $(Y-O)/\mathbb{C}^*$ 在其自然诱导拓扑下不是豪斯多夫空间（Hausdorff space）。

令 $Y_1$ 为 $Y$ 中 $p=0$ 而 $s_i$ 不全为零的子集；令 $Y_2$ 为 $s_i=0$、$p\neq0$ 的子集。任意接近原点的 $\mathbb{C}^*$ 轨道（$O$ 本身除外）恰是 $Y_1$ 或 $Y_2$ 中的轨道。若 $\tilde{Y}=Y-(Y_1\cup Y_2\cup O)$，则由于 $\tilde{Y}$ 只包含“好”的 $\mathbb{C}^*$ 轨道（即闭的、并且与原点保持正距离的轨道），且 $\mathbb{C}^*$ 在 $\tilde{Y}$ 上自由作用，商空间 $\tilde{Y}/\mathbb{C}^*$ 是一个流形，情形就如同 $\mathbb{C}^*$ 是紧群一样。然而，删去 $Y_1$、$Y_2$ 与 $O$ 过于粗糙；适当地纳入某些行为不良的轨道，便可得到 $\tilde{Y}/\mathbb{C}^*$ 的自然部分紧化。

% [Tutor Note: $\tilde{Y}$ 正是 GIT 意义下的稳定点集，$Y_1$、$Y_2$、$O$ 则是随 $r$ 而变的“坏”点（不稳定点，或落在半稳定边界上的点）。按 Mumford 的几何不变量理论，须先选定一个线性化（在这里等价于选定矩映射水平 $r$，即选定 $D$ 函数），再取半稳定点集的商；$r$ 一变，哪些点算半稳定随之改变，于是得到互不相同的 GIT 商——这正是下文 (i)(ii)(iii) 三分支的来源。$\tilde{Y}//U(1)\cong\tilde{Y}/\mathbb{C}^*$ 是 Kempf--Ness 定理的内容：在稳定点集上辛商与全纯商一致。]

在所考虑的简单情形中，不难猜出应当如何做。而且存在一套系统的理论（几何不变量理论（geometric invariant theory）[34]；入门可参见 Newstead 的书 [35]）。与其纯粹从数学上研究这一情形，不如回到第 3.1 节的物理问题。在该问题中，$Y$ 被赋予凯勒度规 $\mathrm{d}\tau^2=\sum_i|\mathrm{d}s_i|^2+|\mathrm{d}p|^2$。此凯勒度规并非 $\mathbb{C}^*$ 不变，故 $\mathbb{C}^*$ 不是第 3 节所研究的 $N=2$ 超对称场论的对称群。但该度规在最大紧子群 $U(1)\subset\mathbb{C}^*$ 下不变（其作用由 (4.1) 给出并取 $|\lambda|=1$）。这个 $U(1)$ 正是第 3 节所用的规范群。与凯勒度规 $\mathrm{d}\tau^2$ 相伴的凯勒形式为 $\omega=-i\sum_i\mathrm{d}\bar{s}_i\wedge\mathrm{d}s_i-i\,\mathrm{d}\bar{p}\wedge\mathrm{d}p$。$Y$ 上的 $U(1)$ 作用保持此凯勒形式。由于 $Y$ 单连通，必然存在一个哈密顿函数，它通过泊松括号生成 $Y$ 上的 $U(1)$ 作用；该函数正是
\begin{equation}
\tilde{D}=-\frac{D}{e^2}=\sum_i|s_i|^2-n|p|^2-r,
\tag{4.2}
\end{equation}
其中 $r$ 为任意常数。因此，在第 3 节中扮演重要角色的 $D$ 函数，正是在此意义下规范群的生成元。

% [Tutor Note: (4.2) 中的 $\tilde{D}$ 就是 $U(1)$ 作用的矩映射。取生成元向量场 $\xi=i\sum_a\big(s_a\partial_{s_a}-\bar{s}_a\partial_{\bar{s}_a}\big)-in\big(p\partial_p-\bar{p}\partial_{\bar{p}}\big)$，在本页 $\omega$ 的约定 $\iota_\xi\omega=-\mathrm{d}\mu$ 下算得 $\mu=\sum_i|s_i|^2-n|p|^2$，即 $\tilde{D}=\mu-r$。于是第 3 节的 $D$ 项方程 $D=0$ 就是“矩映射取常值 $\mu=r$”，而 $r$ 即 FI 参数。易错点：由 (4.2) 有 $D=-e^2(\mu-r)$，故 $D=0$ 与 $\tilde{D}=0$ 等价，但 $D$ 与 $\mu$ 相差一个整体因子 $-e^2$ 与一个常数 $e^2r$，不可直接把两者等同。]

在第 3 节中，取 $D=0$ 是自然的，因为要极小化能量，然后再除以规范群。先令 $D=0$ 再
% ==== 原文 PDF 第 46 页（印刷页 46）====
\label{op:46}
% [跨页接续: 原文第 45 页/第 46 页；本句跨页，此处按原文不切分句子]
对 $U(1)$ 取商，这一复合操作给出所谓 $Y$ 对 $U(1)$ 的辛商。$Y$ 对 $U(1)$ 的辛商常记为 $Y//U(1)$，它当然依赖于 $r$。当 $r$ 经过某些特殊值时 $Y//U(1)$ 奇异，此时甚至连 $Y//U(1)$ 的拓扑都会改变——在我们的例子中这只发生在 $r=0$ 处。而在拓扑不变之时，$Y//U(1)$ 的辛结构仍依赖于 $r$。$Y//U(1)$ 具有自然的辛结构：把 $\omega$ 限制到 $D=0$ 上，再投影到商 $Y//U(1)=\{D=0\}/U(1)$ 即得。这一自然辛结构确定 $H^2(Y//U(1),\mathbb{R})$ 中的一个元素，它线性依赖于 $r$；这与 Duistermaat--Heckman 积分公式 [36] 有关。

% [背景补充: 辛商（symplectic quotient，又称 Marsden--Weinstein--Meyer 约化）即“先取矩映射水平集 $\mu^{-1}(r)$、再对紧群取商”；当群在水平集上自由作用时，商自然继承辛形式。它在代数几何中的对应物是 GIT 商。]

% [Tutor Note: Duistermaat--Heckman 定理断言：辛商上约化辛形式的上同调类 $[\omega_r]\in H^2(Y//U(1),\mathbb{R})$ 随矩映射水平 $r$ 线性变化，其变率由“$D=0$ 上的 $U(1)$ 圆丛”的第一陈类（即该丛曲率的上同调类）给出。本例中，$r>0$ 时 $[\omega_r]$ 是 $\mathbf{CP}^{n-1}$ 超平面类的 $r$ 倍；$r$ 越过 $0$ 变号时商的拓扑本身发生跃变（$r<0$ 的坏点商退化为单点），故 $r=0$ 正是辛结构无法光滑延拓的“跨墙”（wall-crossing）点。]

然而，若考虑的不是 $Y$ 而是 $\tilde{Y}$，则会出现某些简化：$\tilde{Y}//U(1)$ 在某种意义下自然地与 $r$ 无关。为看清这一点，考虑 $D$ 函数在某一 $\mathbb{C}^*$ 轨道上的限制，即
\begin{equation}
\tilde{D}(\lambda)=|\lambda|^2\sum_i|s_i|^2-n|\lambda|^{-2n}|p|^2-r.
\tag{4.3}
\end{equation}
由于身处 $\tilde{Y}$ 之中，$|\lambda|^2$ 与 $|\lambda|^{-2n}$ 的系数均非零。$\tilde{D}(\lambda)$ 是 $|\lambda|$ 的单调函数，当 $\lambda\to\infty$ 时趋于 $+\infty$，当 $\lambda\to0$ 时趋于 $-\infty$。由此，$\tilde{D}(\lambda)=0$ 对 $|\lambda|$ 有唯一解。

% [Tutor Note: 单调性可直接验证：在 $\tilde{Y}$ 上 $A\equiv\sum_i|s_i|^2>0$、$B\equiv|p|^2>0$，故 $\mathrm{d}\tilde{D}/\mathrm{d}|\lambda|=2|\lambda|A+2n^2|\lambda|^{-2n-1}B>0$。关键在于 $\tilde{D}$ 的零点不依赖 $r$ 地给出每条 $\mathbb{C}^*$ 轨道上的唯一代表点，故 $\tilde{Y}//U(1)$ 不随 $r$ 改变；与 $Y//U(1)$ 强烈依赖 $r$ 相对照可知，$r$ 依赖性全部来自“坏点如何被纳入”。]

所给的 $\mathbb{C}^*$ 轨道当然恰好贡献商 $\tilde{Y}/\mathbb{C}^*$ 中的一个点。它也恰好贡献 $\tilde{Y}//U(1)$ 中的一个点，因为要求 $D=0$ 便唯一确定了 $|\lambda|$，而 $\lambda$ 的辐角被 $U(1)$ 作用吸收。因此，$\tilde{Y}//U(1)$ 自然地与 $\tilde{Y}/\mathbb{C}^*$ 重合。

从辛的观点看，如何纳入那些坏点——原点 $O$ 以及 $Y_1$、$Y_2$ 中的点——并无神秘之处。然而，结果依赖于 $r$：

(i) 当 $r>0$ 时，令 $D=0$ 在 $Y_1$ 中可行，而对其余坏点不可行。辛商 $Q=Y_1//U(1)$ 是 $\mathbf{CP}^{n-1}$ 的一个拷贝，其凯勒形式与 $r$ 成正比。于是 $Y//U(1)$ 是 $\tilde{Y}/\mathbb{C}^*$ 与此 $\mathbf{CP}^{n-1}$ 的并。

% ==== 原文 PDF 第 47 页（印刷页 47）====
\label{op:47}
(ii) 当 $r=0$ 时，原点是唯一满足 $D=0$ 的坏点。其辛商是单点，$Y//U(1)$ 是 $\tilde{Y}/\mathbb{C}^*$ 与此点的并。

(iii) 当 $r<0$ 时，令 $D=0$ 在 $Y_2$ 中可行，而对其余坏点不可行。辛商 $Y_2//U(1)$ 是单点，$Y//U(1)$ 是 $\tilde{Y}/\mathbb{C}^*$ 与此点的并。

由于 $Q=Y_1//U(1)$ 与 $Y_1/\mathbb{C}^*$ 相同（两者都是 $\mathbf{CP}^{n-1}$），类似地 $Y_2//U(1)=Y_2/\mathbb{C}^*$（两者都是单点），我们可以纯全纯地重述这些结果。下面就 $r>0$ 与 $r<0$ 两种情形详述：

(i) 当 $r>0$ 时，$Y//U(1)$ 与 $Z=Y'/\mathbb{C}^*$ 相同，其中 $Y'=\tilde{Y}\cup Y_1$ 是 $s_i$ 不全为零的那些点之并。$Z$ 纤维化于 $\mathbf{CP}^{n-1}$ 之上（纤维化的方式是忘掉 $p$）。纤维是复平面的一个拷贝，以 $p$ 为参数。因此 $Z$ 是 $\mathbf{CP}^{n-1}$ 上的复线丛；事实上，由 (4.1) 的变换律可见，$Z$ 是线丛 $\mathcal{O}(-n)$ 的全空间。$Z$ 实际上是一个非紧 Calabi--Yau 流形。为看出这一点，从方程 $\sum_i Q_i=0$ 出发——该方程在第 3 节中曾起重要作用（此处的 $Q_i$ 即 (4.1) 中的指数）。此方程保证 $(n+1)$ 形式 $\Theta=\mathrm{d}s_1\wedge\ldots\wedge\mathrm{d}s_n\wedge\mathrm{d}p$ 是 $\mathbb{C}^*$ 不变的。将 $\Theta$ 与生成 $\mathbb{C}^*$ 作用的向量场缩并，得到处处非零的全纯 $n$ 形式 $\Theta'$，它在 $Y'$ 上的限制是 $Z=Y'/\mathbb{C}^*$ 上某个全纯 $n$ 形式的拉回。

% [Tutor Note: $\Theta$ 在 $\mathbb{C}^*$ 下的权为 $\lambda^{n}\cdot\lambda^{-n}=1$，故不变；而生成元向量场 $v=\sum_i s_i\partial_{s_i}-np\partial_p$ 在 $Y'$ 上处处非零（$v=0$ 只在原点），所以 $\iota_v\Theta$ 是 $Y'$ 上处处非零的全纯 $n$ 形式，且 $\mathbb{C}^*$ 不变，可下降为 $Z$ 上的全纯 $n$ 形式，于是 $Z$ 的典范丛平凡——这就是非紧 Calabi--Yau 的定义。等价地可用陈类核对：$Z$ 为 $\mathcal{O}(-n)$ 的全空间，$K_Z=\pi^*\big(K_{\mathbf{CP}^{n-1}}\otimes\mathcal{O}(n)\big)=\pi^*\big(\mathcal{O}(-n)\otimes\mathcal{O}(n)\big)=\mathcal{O}$，故 $c_1(Z)=0$。]

(iii) 当 $r<0$ 时，$Y//U(1)$ 与 $Y''/\mathbb{C}^*$ 相同，其中 $Y''=\tilde{Y}\cup Y_2$ 是 $p\neq0$ 的区域。于是可用 $(s_i,p)$ 上的 $\mathbb{C}^*$ 作用令 $p=1$。这留下由 $\lambda^n=1$ 定义的 $U(1)$ 子群的残余不变性。除以该群，$r<0$ 时的 $Y//U(1)$ 与 $Z'=\mathbb{C}^n/\mathbb{Z}_n$ 相同。借助 $\sum_i Q_i=0$，这又是一个非紧 Calabi--Yau 流形。

% [Tutor Note: $\lambda^n=1$ 的残余群即 $\mathbb{Z}_n$，它在 $\mathbb{C}^n$ 上的作用为 $(s_1,\ldots,s_n)\mapsto(\lambda s_1,\ldots,\lambda s_n)$；因 $\det=\lambda^n=1$，有 $\mathbb{Z}_n\subset SU(n)$，故保持全纯 $n$ 形式 $\mathrm{d}s_1\wedge\ldots\wedge\mathrm{d}s_n$，典范丛平凡。这一步与 $\sum_i Q_i=0$ 是同一件事的两种说法：权重和为零 $\Leftrightarrow$ 商群含于 $SU(n)$。]

有两个主要结论：

(a) $Y//U(1)$ 对任意 $r$ 都带有自然复结构。这在“物理上”显然必须如此，因为第 3 节所研究模型的低能极限（若将超势取为零）正是以该空间为靶空间的 $N=2$ $\sigma$ 模型。
% ==== 原文 PDF 第 48 页（印刷页 48）====
\label{op:48}
% [跨页接续: 原文第 47 页/第 48 页；第 47 页末为 "…with this target"，第 48 页首词为 "space."，本句跨页，此处按原文不切分句子]
（对 $N=4$ 作类似推理，便得到著名的超凯勒商（hyper-Kähler quotient）[18]。）

(b) $r$ 为正、为负或为零时得到的各个复流形 $Y//U(1)$ 在稠密开集上彼此等价，因为在稠密开集上它们都与 $Y'/\mathbb{C}^*$ 重合。对此的技术术语是：这些复流形双有理等价（birationally equivalent）。事实上，我们还能说得更具体：$Z$ 由胀开 $Z'$ 的原点而得——即以 $\mathbf{CP}^{n-1}\cong Y_1//U(1)$ 的拷贝，替换 $Z'=\mathbb{C}^n/\mathbb{Z}_n$ 中的原点（该原点原本作为 $Y_2//U(1)$ 出现）。

% [背景补充: 双有理等价指两个代数簇（复流形）含有同构的稠密开子集，即存在双有理映射；在复维数 $\ge2$ 时它允许沿子簇收缩/胀开，因此双有理等价的流形拓扑可以不同。]

% [Tutor Note: 上述胀开正是 $\mathbb{C}^n/\mathbb{Z}_n$ 的奇点消解：例外除子 $E\cong\mathbf{CP}^{n-1}$，其法丛为 $\mathcal{O}(-n)|_E$，而 $\mathcal{O}(-n)$ 的全空间恰为 $Z$。$r<0$ 与 $r>0$ 两个相因此由一次胀开（反向为收缩）相联系，这是“极端跃迁”（extremal transition）最简单的例子，也是第 5.5 节时空拓扑改变讨论的原型。]

尽管作者是在一个非常特殊的情形中说明这些想法，这些结论仍是普遍的 [33,34]。给定紧群 $G$ 作用在例如 $Y=\mathbb{C}^n$ 上，或作用在一个射影簇（projective variety）上（并选定了凯勒度规），各种可能的辛商 $Y//G$（由不同的 $D$ 函数选择得到）都带有自然复结构，并且彼此双有理等价。事实上，它们都可以在一个稠密开集上与 $\tilde{Y}/G_{\mathbb{C}}$ 等同，其中 $\tilde{Y}$ 是 $Y$ 中适当的稠密开集，$G_{\mathbb{C}}$ 是 $G$ 的复化（complexification）。

% [Tutor Note: “不同的 $D$ 函数选择”与“复化群 $G_{\mathbb{C}}$ 作用在 $Y$ 上的不同线性化”一一对应：对所有选择，稳定点集上的商都等同于 $\tilde{Y}/G_{\mathbb{C}}$，差别只在于 $D$ 函数决定把哪些边界轨道补进来，这才使各个 $Y//G$ 成为同一双有理类的不同（部分）紧化。]

第 3 节的技巧实际上是一种在靶空间双有理等价的 $\sigma$ 模型之间插值的技巧，其手段是改变 $D$ 函数。该技巧可在没有超势时使用，此时它直接给出某些靶空间双有理等价的 $\sigma$ 模型之间的关系。我们将在第 5.5 节中结合时空拓扑的跃迁重新拾起这一主题。

另一方面，若从 $Y$ 上 $\mathbb{C}^*$ 不变的全纯函数 $W$ 出发——它充当超势——则 $W$ 会下降到任一辛商 $Y//G$ 上，成为全纯函数 $\widehat{W}$，并充当超势。于是，由第 3 节方法得到的、靶空间双有理等价的 $\sigma$ 模型之间的任何关系，都可推广到这些空间带有“共同”超势的情形，即这些超势来自 $Y$ 上一个共同的基本规范不变函数。

% [跨页接续: 原文第 48 页/第 49 页。本句跨页，按 SPEC §2“从句子中间切分时按语义完整取”在此补全整句（第 49 页第 1 行 "holomorphic function on different complex manifolds." 属本句），故 s05a 分片起点应为第 49 页 "For instance, in the case treated in §3…"]
一般而言，双有理等价正是使人们能够以这种方式谈论不同复流形上的“同一”全纯函数所需的条件。

% [Tutor Note: 下降之所以成立，是因为 $W$ 的 $\mathbb{C}^*$ 不变性等价于 $W$ 在每条 $\mathbb{C}^*$ 轨道上为常数，从而在商空间上有定义：$r>0$ 分支上须避开 $Y_2$，$r<0$ 分支上须避开 $Y_1$，但两者在稠密开集上给出同一个 $W$，这正是“共同超势”的几何含义，也是第 5 节推广的基本出发点。]
